\section*{第 \ref{ch:b2:randomvar} 章——离散随机变量}

\begin{solution}{exo:b2:randomvar:1}
\emph{二项\omterm{def:b2:randomvar:law}{分布}：} 令 $X = \sum_{i=1}^n X_i$，其中各
$X_i$ 为独立的伯努利变量；$V(X_i) = \E(X_i^2) - \E(X_i)^2 = p - p^2$，
且独立变量的\omterm{def:b2:randomvar:variance}{方差}相加
（\cref{thm:b2:randomvar:variancerules}）：
\[
\E(X) = np, \qquad V(X) = np(1-p) .
\]
\emph{泊松\omterm{def:b2:randomvar:law}{分布}：} $\E\bigl(X(X-1)\bigr) =
\sum_{k\geq2}k(k-1)e^{-\lambda}\frac{\lambda^k}{k!} = \lambda^2
e^{-\lambda}\sum_{j\geq0}\frac{\lambda^j}{j!} = \lambda^2$，于是
\[
V(X) = \E(X^2) - \E(X)^2
= \lambda^2 + \lambda - \lambda^2 = \lambda .
\]
\emph{几何\omterm{def:b2:randomvar:law}{分布}}（$q = 1 - p$）：对
$\sum_{k\geq0}q^k = \frac{1}{1-q}$ 在收敛圆内求导两次
（\cref{ch:b2:powerseries}），$\sum_{k\geq2}k(k-1)q^{k-2} =
\frac{2}{(1-q)^3}$，于是
\[
\E\bigl(X(X-1)\bigr) = pq\sum_{k\geq2}k(k-1)q^{k-2}
= \frac{2q}{p^2},
\qquad
V(X) = \frac{2q}{p^2} + \frac1p - \frac{1}{p^2}
= \frac{q}{p^2} = \frac{1-p}{p^2} .
\]
\end{solution}

\begin{solution}{exo:b2:randomvar:2}
\emph{和：}对 $n \in \N$，由\omterm{def:b2:proba:space}{事件}的不交性及独立性，
\[
\P(X + Y = n)
= \sum_{k=0}^n \P(X = k)\P(Y = n - k)
= e^{-(\lambda + \mu)}\frac{1}{n!}
\sum_{k=0}^n \binom nk \lambda^k\mu^{n-k}
= e^{-(\lambda+\mu)}\frac{(\lambda + \mu)^n}{n!}
\]
由二项式定理，$X + Y \sim \mathcal{P}(\lambda + \mu)$。
\emph{条件\omterm{def:b2:randomvar:law}{分布}：}对 $0 \leq k \leq n$，
\[
\P(X = k \mid X + Y = n)
= \frac{\P(X = k)\P(Y = n - k)}{\P(X + Y = n)}
= \binom nk
\Bigl(\frac{\lambda}{\lambda+\mu}\Bigr)^{k}
\Bigl(\frac{\mu}{\lambda+\mu}\Bigr)^{n-k} ,
\]
即二项\omterm{def:b2:randomvar:law}{分布} $\mathcal{B}\bigl(n,
\frac{\lambda}{\lambda+\mu}\bigr)$：已知总数时，每个
\omterm{def:b2:proba:space}{事件}独立地以
与其速率成正比的概率“选择”第一个来源。
\end{solution}

\begin{solution}{exo:b2:randomvar:3}
当还缺 $k$ 种玩具时，每个新盒子
以概率 $\frac kn$ 带来一种新玩具，与过去独立：
下一种新玩具的等待时间 $W_k$ 服从几何\omterm{def:b2:randomvar:law}{分布}，
$\mathcal{G}\bigl(\frac kn\bigr)$，其中 $\E(W_k) = \frac nk$，且
$T_n = W_n + W_{n-1} + \dots + W_1$（第一个盒子必定带来新玩具，
即 $W_n = 1$，与 $\E = n/n$ 相符）。由线性性，
\[
\E(T_n) = \sum_{k=1}^n \frac nk = n\sum_{k=1}^n\frac1k
\sim n\ln n ,
\]
利用 $\sum_{k\leq n}\frac1k = \ln n + \gamma + o(1)$
（\cref{ch:b2:comparison}）。收集最后几种玩具，才是
真正的代价：一半的盒子都要用在最后这几种上。
\end{solution}

\begin{solution}{exo:b2:randomvar:4}
逐点地，$X(\omega) = \#\{n \geq 1 : X(\omega) \geq n\} =
\sum_{n\geq1}\mathbf{1}_{X \geq n}(\omega)$。二重族
$\bigl(\mathbf{1}_{X \geq n}(\omega)\,\P(\{\omega\})\bigr)_{n,
\omega}$ 非负，因此族的富比尼定理
（\cref{ch:b2:series}）无条件适用：先对
$n$ 求和得 $\E(X)$，先对 $\omega$ 求和得 $\sum_n \P(X
\geq n)$；二者同时有限且相等。对 $X \sim
\mathcal{G}(p)$：$\P(X \geq n) = q^{n-1}$（$q = 1-p$），于是 $\E(X) =
\sum_{n\geq1}q^{n-1} = \frac{1}{1 - q} = \frac1p$.
\end{solution}

\begin{solution}{exo:b2:randomvar:5}
\emph{对称性：}第 $i$ 个被抽出的球是
罐中均匀随机选出的一个球（$N$ 个球中每一个都等可能地落在
抽取顺序的第 $i$ 位上），因此 $\P(Y_i = 1) = \frac MN =
p$，由线性性得 $\E(X) = np$，无需独立性。

\emph{\omterm{def:b2:randomvar:variance}{协方差}：}对 $i \neq j$，$\E(Y_iY_j) = \P(\text{第 }
i, j \text{ 次均为白球}) = \frac{M(M-1)}{N(N-1)}$（两个不同位置的有序对被均匀地赋予
两个不同球的有序对）。
由此
\[
\operatorname{Cov}(Y_i, Y_j)
= \frac{M(M-1)}{N(N-1)} - \frac{M^2}{N^2}
= \frac{M(N - M)}{N^2}\cdot\frac{-1}{N-1}
= -\frac{p(1-p)}{N-1} < 0 :
\]
抽到一个白球，会使其余各次抽取中的白球变得更少。
由\cref{thm:b2:randomvar:variancerules}，
\[
V(X) = np(1-p) + n(n-1)\Bigl(-\frac{p(1-p)}{N-1}\Bigr)
= np(1-p)\,\frac{N - n}{N - 1} \leq np(1-p) :
\]
不放回抽样的均值相同，但\omterm{def:b2:randomvar:variance}{方差}\emph{更小}
比放回抽样时（仅当 $n = 1$ 时取等号），
负相关起到了稳定作用。当 $n = N$ 时
\omterm{def:b2:randomvar:variance}{方差}为零，此时计数是确定的。
\end{solution}

\begin{solution}{exo:b2:randomvar:6}
围绕 $m = \E(X)$ 展开：
\[
\E\bigl((X - c)^2\bigr)
= \E\bigl((X - m)^2\bigr) + 2(m - c)\,\E(X - m) + (m - c)^2
= V(X) + (m - c)^2 ,
\]
仅在 $c = m$ 时取最小值 $V(X)$——\omterm{def:b2:randomvar:expectation}{期望}是
均方意义下最好的常数预测。

若 $\P(X = m) = 1$，则 $(X - m)^2$ 以概率
$1$ 为零，故 $V(X) = 0$（定义式的各项在
一个零测集之外为零）。反之，若 $V(X) = 0$，由切比雪夫不等式
（\cref{thm:b2:randomvar:markov}）给出 $\P\bigl(\abs{X - m} \geq
\frac1n\bigr) \leq n^2\,V(X) = 0$ 对每个 $n$ 成立；\omterm{def:b2:proba:space}{事件}
$\bigl\{\abs{X - m} \geq \frac1n\bigr\}$ 单调递增趋于 $\{X \neq
m\}$，由单调\omterm{def:b2:metric:continuity}{连续}性（\cref{thm:b2:proba:continuity}）
可得 $\P(X \neq m) = 0$。
\end{solution}

\begin{solution}{exo:b2:randomvar:7}
$\E(S_n) = \frac n2$，$V(S_n) = \frac n4$。
\emph{马尔可夫：}$\P\bigl(S_n \geq \frac{3n}4\bigr) \leq
\frac{n/2}{3n/4} = \frac23$——这是一个常数界，对
较大的 $n$ 毫无用处。
\emph{切比雪夫：}\omterm{def:b2:proba:space}{事件}蕴含 $\abs{S_n - \frac n2} \geq
\frac n4$，故概率 $\leq \frac{n/4}{(n/4)^2} =
\frac4n$——趋于零，但只是多项式衰减。
\emph{切尔诺夫：}由独立性，$\E(e^{tS_n}) =
\prod_{i=1}^n\E(e^{tX_i}) = \bigl(\frac{1 + e^t}{2}\bigr)^n$，且
对 $e^{tS_n} \geq e^{3nt/4}$ 应用马尔可夫不等式，对每个 $t >
0$，有
\[
\P\Bigl(S_n \geq \frac{3n}4\Bigr)
\leq \Bigl(\frac{1 + e^t}{2}\Bigr)^n e^{-3nt/4}
= \exp\Bigl(n\bigl(\ln\tfrac{1 + e^t}{2} - \tfrac{3t}4\bigr)\Bigr).
\]
求指数的最小值：$\frac{\dd}{\dd t}\ln\frac{1+e^t}{2} =
\frac{e^t}{1 + e^t} = \frac34$ 当 $e^t = 3$ 时成立，即 $t = \ln 3$，
于是得到
\[
\P\Bigl(S_n \geq \frac{3n}4\Bigr)
\leq \Bigl(\frac{4}{2}\Bigr)^n 3^{-3n/4}
= \bigl(2 \cdot 3^{-3/4}\bigr)^n \approx (0.877)^n ,
\]
指数级地小。这一由马尔可夫 $\to$ 切比雪夫 $\to$
再到切尔诺夫的层级是标准的阶梯：每一级都
在变量的一个增长更快的函数上应用马尔可夫不等式。
\end{solution}

\begin{solution}{exo:b2:randomvar:8}
由\cref{thm:b2:randomvar:transfer}的变换公式
应用于 $f\bigl(\frac{S_n}{n}\bigr)$，其中 $S_n \sim
\mathcal{B}(n, x)$:
\[
\E\Bigl[f\Bigl(\frac{S_n}{n}\Bigr)\Bigr]
= \sum_{k=0}^n f\Bigl(\frac kn\Bigr)\binom nk x^k(1-x)^{n-k}
= B_nf(x) .
\]
取 $\delta > 0$，并把 $\abs{f(S_n/n) - f(x)}$ 在\omterm{def:b2:proba:space}{事件} $D
= \bigl\{\abs{\frac{S_n}{n} - x} \geq \delta\bigr\}$ 上拆开：
在 $D$ 之外，差值至多为\omterm{def:b2:metric:continuity}{连续}模
$\omega_f(\delta) = \sup_{\abs{s - t}\leq\delta}\abs{f(s) -
f(t)}$；在 $D$ 上，至多为 $2\norm f_\infty$。取\omterm{def:b2:randomvar:expectation}{期望}，并
利用切比雪夫不等式，其中 $V\bigl(\frac{S_n}{n}\bigr) =
\frac{x(1-x)}{n} \leq \frac{1}{4n}$：
\[
\abs{B_nf(x) - f(x)}
\leq \E\,\abs{f(S_n/n) - f(x)}
\leq \omega_f(\delta)
+ 2\norm f_\infty\,\P(D)
\leq \omega_f(\delta) + \frac{2\norm f_\infty}{4n\delta^2} .
\]
由 $f$ 在 $[0, 1]$ 上的一致\omterm{def:b2:metric:continuity}{连续}性，使得 $\omega_f(\delta) \to
0$：先取 $\delta$，再取 $n$，于是 $B_nf \to f$ \omterm{def:b2:funcseq:def}{一致收敛}——这正是
在\cref{ch:b2:funcseq}中的魏尔斯特拉斯逼近定理，其
“计数引理”现在可被识别为
二项\omterm{def:b2:randomvar:law}{分布}律上的切比雪夫不等式。
\end{solution}

\begin{solution}{exo:b2:randomvar:9}
展开 $S_n^4 = \sum_{i,j,k,l}X_iX_jX_kX_l$，并取
\omterm{def:b2:randomvar:expectation}{期望}。由独立性与中心化，任何含有
恰出现一次的下标的项为零（$\E(X_i) = 0$ 的因子
可提出）。保留下来的项有：$n$ 个对角项 $\E(X_i^4)$，以及
下标两两相等的项，$\E(X_i^2X_j^2) =
\E(X_1^2)^2$，其中 $i \neq j$，共出现 $3n(n-1)$ 次：先选
无序的一对取值（$\binom n2$ 种方式），然后
$\frac{4!}{2!\,2!} = 6$ 种方式把它们放入四个位置——
$6\binom n2 = 3n(n-1)$。于是，由 $\E(X_1^2)^2
\leq \E(X_1^4)$（由琴生不等式或柯西–施瓦茨不等式），
\[
\E(S_n^4) = n\,\E(X_1^4) + 3n(n-1)\,\E(X_1^2)^2
\leq C n^2,
\qquad C = 4\,\E(X_1^4) .
\]
4 阶马尔可夫不等式：
\[
\P\Bigl(\Bigl|\frac{S_n}{n}\Bigr| \geq \varepsilon\Bigr)
= \P\bigl(S_n^4 \geq n^4\varepsilon^4\bigr)
\leq \frac{Cn^2}{n^4\varepsilon^4}
= \frac{C}{n^2\varepsilon^4} ,
\]
这是一个收敛级数。由博雷尔–坎泰利第一引理
（\cref{thm:b2:proba:borelcantelli}），对每个 $j$，\omterm{def:b2:proba:space}{事件}
$B_j = \limsup_n\bigl\{\abs{S_n/n} \geq \frac1j\bigr\}$ 的概率
为 $0$，故 $\P\bigl(\bigcup_j B_j\bigr) = 0$（由
可数次可加性）。在补集上——其概率为 $1$
——对每个 $j$，存在 $N$，使得 $\abs{S_n/n} < \frac1j$ 对
所有 $n \geq N$ 成立，即 $S_n/n \to 0$。强大数定律
在四阶矩假设下成立；去掉此假设
（柯尔莫哥洛夫定理）属于本科三年级的内容。
\end{solution}

\begin{solution}{exo:b2:randomvar:10}
$\P(M \leq k) = \bigl(\frac k6\bigr)^2$（两颗骰子都不超过
$k$，且相互独立），故 $\P(M \geq k) = 1 -
\bigl(\frac{k-1}6\bigr)^2$，且
\[
\E(M) = \sum_{k=1}^6\P(M \geq k)
= 6 - \frac{0 + 1 + 4 + 9 + 16 + 25}{36}
= 6 - \frac{55}{36} = \frac{161}{36} \approx 4.47 ,
\]
明显高于单颗骰子的均值 $3.5$，
这正是最大值应有的情形。
\end{solution}

\begin{solution}{exo:b2:randomvar:11}
设 $I_i = \mathbf 1_{\sigma(i) = i}$，则 $\P(\sigma(i) = i) =
\frac{(n-1)!}{n!} = \frac1n$，故 $\E(F_n) = n\cdot\frac1n =
1$。当 $i \neq j$ 时，$\P(\sigma(i) = i, \sigma(j) = j) =
\frac{(n-2)!}{n!} = \frac1{n(n-1)}$，因此
\[
\operatorname{Cov}(I_i, I_j) = \frac1{n(n-1)} - \frac1{n^2}
= \frac{1}{n^2(n-1)} .
\]
由\omterm{def:b2:randomvar:variance}{方差}的运算规则
（\cref{thm:b2:randomvar:variancerules}），
\[
V(F_n) = n\cdot\frac1n\Bigl(1 - \frac1n\Bigr)
+ n(n-1)\cdot\frac1{n^2(n-1)}
= 1 - \frac1n + \frac1n = 1 .
\]
均值为 $1$，\omterm{def:b2:randomvar:variance}{方差}为 $1$，与 $n$ 无关——与
匹配问题的泊松极限一致
（见\cref{exo:b2:proba:5}）。
\end{solution}

\begin{solution}{exo:b2:randomvar:12}
设 $T_n = \sum_{k=1}^nG_k$，其中 $G_k \sim \mathcal G(k/n)$ 为
缺少 $k$ 个时见到新玩具所需的时间，各阶段
相互独立。由此
\[
V(T_n) = \sum_{k=1}^n\frac{1 - k/n}{(k/n)^2}
\leq \sum_{k=1}^n\frac{n^2}{k^2}
\leq \frac{\pi^2}6\,n^2 ,
\]
由\cref{ex:b2:fourier:basel}可得。又 $\E(T_n) = nH_n$，$H_n =
\sum_1^n\frac1k$（见\cref{exo:b2:randomvar:3}），由切比雪夫不等式得，
对 $\varepsilon > 0$，
\[
\P\bigl(\abs{T_n - nH_n} \geq \varepsilon\,n\ln n\bigr)
\leq \frac{\pi^2n^2/6}{\varepsilon^2n^2\ln^2 n}
= \frac{\pi^2}{6\,\varepsilon^2\ln^2n}
\xrightarrow[n\to\infty]{} 0 .
\]
由于 $H_n \sim \ln n$，除以 $n\ln n$ 可知
$T_n/(n\ln n) \to 1$ 依概率成立：$T_n$ 的波动
量级为 $n$，与均值 $n\ln n$ 相比可以忽略。
\end{solution}

\begin{solution}{pb:b2:randomvar:1}
\textbf{1.} $\E(S_n) = \frac n2$，由马尔可夫不等式
（\cref{thm:b2:randomvar:markov}）可得 $\P(S_n \geq an) \leq
\frac{n/2}{an} = \frac1{2a}$。马尔可夫不等式只利用均值：它
无法区分集中在 $n/2$ 处的\omterm{def:b2:randomvar:law}{随机变量}与
\omterm{def:b2:randomvar:law}{分布}在 $0$ 与 $n$ 之间的变量，故它对尾概率的估计
就好像全部质量都可能落在那里。

\textbf{2.} 公平的二项\omterm{def:b2:randomvar:law}{分布}关于 $n/2$ 对称
（$S_n$ 与 $n - S_n$ 同\omterm{def:b2:randomvar:law}{分布}），故取 $x = n(a -
\frac12) > 0$，则两\omterm{def:b2:proba:space}{事件} $\{S_n - \frac n2 \geq x\}$ 与
$\{S_n - \frac n2 \leq -x\}$ 互不相交且概率相等：
$\P(S_n \geq an) = \frac12\P(\abs{S_n - \frac n2} \geq x)$。
由切比雪夫不等式，$V(S_n) = \frac n4$：
\[
\P(S_n \geq an)
\leq \frac12\cdot\frac{n/4}{n^2(a - 1/2)^2}
= \frac1{8n(a - 1/2)^2},
\]
它等于 $\frac2n$，此时 $a = \frac34$。

\textbf{3.} 由独立性及乘积定理，
$\E(\eu^{tS_n}) = \bigl(\E \eu^{tX_1}\bigr)^n = \bigl(\frac{1
+ \eu^t}2\bigr)^n$。对 $\eu^{tS_n}$ 应用马尔可夫不等式：
\[
\P(S_n \geq an) \leq \eu^{-tan}\Bigl(\frac{1 +
\eu^t}2\Bigr)^{\!n} = \exp\Bigl(n\bigl(\ln\tfrac{1 +
\eu^t}2 - ta\bigr)\Bigr).
\]
指数关于 $t$ 的导数为 $\frac{\eu^t}{1 + \eu^t}
- a$，在 $\eu^t = \frac a{1-a}$ 处为零，即 $t^* =
\ln\frac a{1-a} > 0$；此时 $\frac{1 + \eu^{t^*}}2 =
\frac1{2(1-a)}$，指数等于
\[
n\Bigl(-\ln 2 - \ln(1-a) - a\ln\frac a{1-a}\Bigr)
= -n\bigl(\ln2 + a\ln a + (1-a)\ln(1-a)\bigr) = -n\,I(a),
\]
其中 $I(\frac12) = 0$，且 $I'(a) = \ln\frac a{1-a} > 0$
在 $\intoo{\frac12}1$ 上，故 $I(a) > 0$。当 $a = \frac34$ 时：
$\eu^{-I(3/4)} = \frac12(\tfrac34)^{-3/4}(\tfrac14)^{-1/4} =
2\cdot3^{-3/4}$，即\cref{exo:b2:randomvar:7}给出的界。

\textbf{4.} 这 $n + 1$ 个数 $\binom nja^j(1-a)^{n-j}$ 之和
为 $1$，最大者是 $j = k = an$ 处的那一个（$\mathcal B(n, a)$
的众数为 $\floor{(n+1)a} = k$，此处成立）。
$n + 1$ 个和为 $1$ 的数，其最大者至少为
$\frac1{n+1}$：
\[
\binom nk a^k(1-a)^{n-k} \geq \frac1{n+1}
\quad\Longrightarrow\quad
\binom nk \geq \frac{a^{-an}(1-a)^{-n(1-a)}}{n+1}
= \frac{\eu^{nH(a)}}{n+1}.
\]
由此 $\P(S_n \geq an) \geq \binom{n}{an}2^{-n} \geq
\eu^{n(H(a) - \ln2)}/(n+1) = \eu^{-nI(a)}/(n+1)$：在
相差多项式因子 $n + 1$ 的意义下，切尔诺夫指数即为真值。

\textbf{5.} 取 $n = 100$，$a = \frac34$：马尔可夫界为 $\frac23$；
切比雪夫界为 $\frac2{100} = 0.02$；
切尔诺夫界为 $(2\cdot3^{-3/4})^{100} = \eu^{-100\,I(3/4)} \approx
2.1\cdot10^{-6}$，与精确值 $2.8\cdot10^{-7}$ 相比。寓意：
每多利用一阶矩信息，界只按多项式因子缩小；
指数矩则改变了界的\emph{性质}。

\textbf{6.} $\cosh t = \sum_{k\geq0}\frac{t^{2k}}{(2k)!}$ 和
$\eu^{t^2/2} = \sum_{k\geq0}\frac{t^{2k}}{2^kk!}$；由逐项比较即得所要的结论，
只需 $(2k)! \geq 2^kk!$，该式由
归纳法成立：$(2k)! = 2k(2k-1)\cdot(2k-2)! \geq 2k\cdot
2^{k-1}(k-1)! = 2^kk!\cdot(2k-1) \geq 2^kk!$。

\textbf{7.} 由独立性，$\E\bigl(\eu^{t\sum\varepsilon_i}
\bigr) = (\cosh t)^n \leq \eu^{nt^2/2}$，于是由马尔可夫不等式得
$\P(\sum\varepsilon_i \geq s) \leq \eu^{nt^2/2 - ts}$；
使其在 $t = s/n$ 处最小，得 $\eu^{-s^2/(2n)}$。

\textbf{8.} 取 $X_i = \frac{1 + \varepsilon_i}2$，$\widehat
p_n - \frac12 = \frac1{2n}\sum\varepsilon_i$，于是
$\{\widehat p_n - \frac12 \geq \delta\} =
\{\sum\varepsilon_i \geq 2n\delta\}$，由第 7 题得到
界 $\eu^{-(2n\delta)^2/(2n)} = \eu^{-2n\delta^2}$。对称\omterm{def:b2:proba:space}{事件}的界相同，
由此得到因子 $2$，用于
$\abs{\widehat p_n - \frac12} \geq \delta$。

\textbf{9.} $\E(\eu^{tX}) = \sum_x\eu^{tx}\P(X = x)$ 是
关于 $t$ 的光滑函数的级数，其逐项导数
在每个紧的 $t$ 区间上，均被
$\eu^{\abs t}\P(X = x)$（因 $0 \leq x \leq 1$）所控制。由
\omterm{def:b2:funcseq:series}{正规收敛}级数的求导定理
（\cref{thm:b2:funcseq:differentiation}）知其
两次可导，并由商法则得 $\psi' = \E_t(X)$
及 $\psi'' = \E_t(X^2) - \E_t(X)^2$，其中 $\E_t$ 为
对重新加权后的权重的\omterm{def:b2:randomvar:expectation}{期望}，
$\eu^{tx}\P(X{=}x)/\E(\eu^{tX})$——非负，和为
$1$，支撑仍为相同的取值 $x \in \intcc01$。取值于
$\intcc01$ 的变量的\omterm{def:b2:randomvar:variance}{方差}至多为 $\frac14$：由
\cref{exo:b2:randomvar:6}，它等于 $\min_c\E_t((X - c)^2) \leq
\E_t\bigl((X - \tfrac12)^2\bigr) \leq \tfrac14$。带
积分型余项的泰勒公式，利用 $\psi(0) = 0$，$\psi'(0) = p$：
\[
\psi(t) = tp + \int_0^t(t - s)\,\psi''(s)\,\dd s
\leq tp + \frac{t^2}2\cdot\frac14,
\]
即 $\E(\eu^{t(X - p)}) \leq \eu^{t^2/8}$，对一切实数 $t$ 成立。

\textbf{10.} 由独立性，$\E\bigl(\eu^{t(S_n -
np)}\bigr) \leq \eu^{nt^2/8}$；由马尔可夫不等式并取 $t
= 4\delta$ 得
\[
\P(\widehat p_n - p \geq \delta)
\leq \eu^{nt^2/8 - tn\delta}\Big|_{t = 4\delta}
= \eu^{-2n\delta^2};
\]
将此结论用于变量 $1 - X_i$（它们同样取值于
$\intcc01$）可界定另一侧尾概率，因而双侧界为
$2\eu^{-2n\delta^2}$。

\textbf{11.} 切比雪夫不等式只需二阶矩，给出
$\frac{p(1-p)}{n\delta^2}$；霍夫丁不等式需要
\emph{有界性}，给出 $2\eu^{-2n\delta^2}$。当 $p =
\frac12$，$\delta = 0.03$ 时界为：
$\frac{278}{n}$（近似）与
$2\eu^{-0.0018n}$；二者约在 $n \approx 1200$ 处相交，此后
指数界占优，而且优势极为明显（$n = 5000$：
$0.056$ 与 $2.5\cdot10^{-4}$）。

\textbf{12.} 由霍夫丁不等式（第 10 题），$\P(\abs{\widehat
p_n - p} \geq \delta) \leq 2\eu^{-2n\delta^2} \leq \alpha$ 成立，
只要 $2n\delta^2 \geq \ln\frac2\alpha$，即 $n \geq
\frac{\ln(2/\alpha)}{2\delta^2}$。

\textbf{13.} 取 $\delta = 0.03$，$\alpha = 0.05$：$n \geq
\frac{\ln 40}{2\cdot0.0009} \approx 2049.4$：需 $2050$ 人。
当 $\delta = 0.01$ 时：$n \geq \frac{\ln40}{0.0002} \approx
18\,445$。总体规模从不出现，因为每位被抽到的
选民都被建模为一次新的伯努利$(p)$ 抽样：民意调查的
难度是一枚硬币的\omterm{def:b2:randomvar:variance}{方差}，而非
国家的大小。将误差幅度减半，样本量要增至四倍
——即 $1/\delta^2$ 规律。

\textbf{14.} 切比雪夫不等式：$\P(\abs{\widehat p_n - p} \geq
\delta) \leq \frac{p(1-p)}{n\delta^2} \leq
\frac1{4n\delta^2} \leq \alpha$，只要 $n \geq
\frac1{4\alpha\delta^2}$，即三点处为 $5556$——
约为霍夫丁要求的 $2.7$ 倍。不放回时
\omterm{def:b2:randomvar:variance}{方差}要乘以 $\frac{N -
n}{N-1} < 1$（\cref{exo:b2:randomvar:5}），故同样的 $n$ 只会
更好：放回情形的计算是
偏于保守的那一个。

\textbf{15.} $\{\widehat p_n \leq \frac12\} \subseteq
\{\widehat p_n - 0.52 \leq -0.02\}$，于是由单侧
霍夫丁不等式得 $\P(\widehat p_n \leq \tfrac12) \leq
\eu^{-2n(0.02)^2} \leq 0.01$，只要 $n \geq \frac{\ln
100}{2\cdot0.0004} \approx 5756.5$：需 $5757$ 名选民。代价
与\emph{领先幅度}的平方成反比，而不是
由所要求的精度决定：势均力敌的选举代价高昂。

\textbf{16.} 所用的假设为：样本从选民中均匀抽取，
且各次抽取相互独立；每位被抽到的人都
诚实作答，且 $p$ 在调查期间不变。真实的
民意调查违反了上述三条：可接触且愿意作答的受访者
并非均匀样本（存在选择偏差与无应答偏差），
而回答可能不真实或不稳定。这些属于
\emph{偏差}误差：它们使 $\E(\widehat p_n)$ 偏离
$p$ 的量与 $n$ 无关，故任何样本量都
无法减小它们——本部分的数学只控制
波动项。

\textbf{17.} 对大小为 $m$ 的一组，用切比雪夫不等式：$\P(\abs{
\widehat p^{(i)} - p} \geq \delta) \leq \frac{1}{4m\delta^2}
\leq \frac18$，当 $m \geq \frac2{\delta^2}$ 时成立。若少于
$k/2$ 组出错，则多于 $k/2$ 个取值
$\widehat p^{(i)}$ 落在开区间 $\intoo{p -
\delta}{p + \delta}$ 内，它们的中位数亦然；由此
$\{\abs{M - p} \geq \delta\}$ 蕴含至少 $\lceil
k/2\rceil$ 组出错（共 $k$ 个独立分组）。
对 $\binom k{\lceil k/2\rceil}$ 种可能的出错分组集合取并界，
得到：
\[
\P(\abs{M - p} \geq \delta)
\leq \binom{k}{\lceil k/2\rceil}
\Bigl(\frac18\Bigr)^{k/2}
\leq 2^k\,8^{-k/2} = 2^{-k/2} :
\]
关于分组数的指数衰减，仅凭
\omterm{def:b2:randomvar:variance}{方差}即可获得——正适用于求和项无界、
霍夫丁不等式不可用的情形。

\textbf{18.} 柯西–施瓦茨不等式
（见\cref{thm:b2:randomvar:jensen}）：
\[
\E(X) = \E(X\,\mathbf 1_{X>0})
\leq \sqrt{\E(X^2)}\sqrt{\E(\mathbf 1_{X>0}^2)}
= \sqrt{\E(X^2)\,\P(X > 0)} ;
\]
平方后相除。

\textbf{19.} 设 $h(a) = I(a) - 2(a - \tfrac12)^2$，则
$h(\tfrac12) = 0$, $h'(a) = \ln\frac a{1-a} - 4(a -
\tfrac12)$ 在 $\tfrac12$ 处为零，且
\[
h''(a) = \frac1a + \frac1{1-a} - 4 = \frac{1}{a(1-a)} - 4
\geq 0
\]
由 $a(1-a) \leq \frac14$ 可知，$h'$ 自 $0$ 起在
$\intco{\frac12}1$ 上递增，故 $h' \geq 0$，从而 $h \geq 0$：$I(a)
\geq 2(a - \tfrac12)^2$。

\textbf{20.} $I(\tfrac12) = I'(\tfrac12) = 0$, $I''(a) =
\frac1{a(1-a)}$ 给出 $I''(\tfrac12) = 4$，且 $I'''(\tfrac12)
= 0$（函数关于 $\tfrac12$ 对称），故
$I(\tfrac12 + \delta) = 2\delta^2 + O(\delta^4)$。第 4 问随后
给出真实尾概率的\emph{下界}
$\eu^{-n(2\delta^2 + O(\delta^4))}/(n+1)$：对于小的 $\delta$，
霍夫丁指数 $2n\delta^2$ 渐近精确，
仅可能有关于 $n$ 的多项式量级的改进。

\textbf{21.} 马尔可夫不等式：一阶矩，衰减为 $1/a$，仅
作为其余各界的核心（第 1 问显示其平坦）。
切比雪夫不等式：二阶矩，衰减为 $\frac{V}{n\delta^2}$，
无需进一步假设即为精确的
（见\cref{ex:b2:randomvar:chebsharp}），并且是第 23 问的最佳工具。
四阶矩（见\cref{exo:b2:randomvar:9}）：
衰减为 $C/n^2$，恰好足以给出强大数律。
霍夫丁不等式：有界变量，衰减为 $2\eu^{-2n\delta^2}$，
是第三部分的主力工具。切尔诺夫界取精确速率
$I(a)$：拥有完整的指数矩，指数不可超越
（第 4、20 问），是其余一切界的参照点。

\textbf{22.} 若硬币公平：$\P(\widehat p_n > 0.525)
\leq \P(\widehat p_n - \tfrac12 \geq 0.025) \leq
\eu^{-2n(0.025)^2}$。若 $p = 0.55$：$\P(\widehat p_n \leq
0.525) \leq \P(\widehat p_n - 0.55 \leq -0.025) \leq
\eu^{-2n(0.025)^2}$。两类错误概率均低于 $0.01$，当
$2n(0.025)^2 \geq \ln 100$，即 $n \geq 3684.2$：$3685$
次抛掷。（区分相距 $2.5$ 个百分点的两个假设，
其代价与估计到 $\pm2.5$ 个百分点时相同。）

\textbf{23.} 霍夫丁不等式：$n \geq \frac{\ln 40}{2(0.005)^2}
\approx 73\,778$。切比雪夫不等式取真实\omterm{def:b2:randomvar:variance}{方差} $p(1-p)
= 0.0099$: $n \geq \frac{0.0099}{0.05\cdot(0.005)^2} =
7920$——所需样本量少九倍。霍夫丁的指数
$2n\delta^2$ 把\omterm{def:b2:randomvar:variance}{方差}取为最坏情形 $\frac14$，
在 $p = 0.01$ 时荒谬地悲观；朴素的二阶矩
则更为明智。缺少的工具是考虑\omterm{def:b2:randomvar:variance}{方差}的
指数型界（伯恩斯坦不等式，本科三年级分册）——或
对于稀有\omterm{def:b2:proba:space}{事件}，可用在
\cref{ch:b2:genfun}中所证的泊松近似，它在自然的相对
尺度上奏效。

\textbf{24.} 固定 $\delta > 0$：$\sum_n 2\eu^{-2n\delta^2} <
\infty$（几何型级数），故博雷尔–坎泰利~1
（见\cref{thm:b2:proba:borelcantelli}）给出
$\P(\abs{\widehat p_n - p} \geq \delta \text{ 无穷多次
}) = 0$，即\omterm{def:b2:proba:space}{事件} $E_j =
\bigcup_N\bigcap_{n\geq N}\{\abs{\widehat p_n - p} <
\tfrac1j\}$ 的概率为 $1$，对每个 $j$ 成立。可数
交 $\bigcap_jE_j$ 的概率仍为 $1$，
（由补集上的次可加性），且在其上 $\widehat p_n
\to p$：抛掷硬币的强大数律，其中
有界性所扮演的角色，正如四阶矩在
\cref{exo:b2:randomvar:9}中所扮演的那样。

\textbf{25.} 一阶矩只得到平坦的界；二阶矩得到
$1/(n\delta^2)$，不能更多（即精确性例子）；四阶矩
得到 $1/n^2$，足以通过裂项求和得到几乎必然成立的定律；
有界性得到 $\eu^{-2n\delta^2}$；而完整的
指数矩得到精确速率 $I$，任何方法都
无法超越。调查 $2050$ 人对任何国家都足够，因为
样本的波动由硬币的\omterm{def:b2:randomvar:variance}{方差}决定，
而非总体的大小——$1/\delta^2$ 与
$\ln(1/\alpha)$ 的代价标签是普适的。本科三年级分册的
中心极限定理在
$\sqrt n$ 尺度上，以精确极限律取代这些不等式，
并给出显式常数——使本题的每个界化为
渐近等式。
\end{solution}
